%!TEX root = ../main.tex
% =============================================================================
%  二、模型假设与符号说明        建议 2 页
%  评审在问：你的假设合理吗？符号清楚吗？
%  本节必须出现：① 假设表 A1–A8，每条含「必要性 + 影响方向 + 验证方式」
%                ② 完整符号表，含单位，且\textbf{必须含对偶变量}
%                ③ 信息结构 R1/R2/R3 的显式声明
%  判据：一条假设若不能回答"不用它模型会怎样"，它就不该出现在假设表里。
% =============================================================================

\section{模型假设与符号说明}

\subsection{模型假设}

题目已给定储能设备的最大容量 12000~kWh、最大充放电功率 5000~kW、
2025 年 1 月 1 日 \(0{:}00\) 储电量 6000~kWh、运行区间 1200--10800~kWh 与充放电效率 90\%。
除此之外，为使模型可解且与现实相符，本文作如下假设。
每条假设均给出\textbf{必要性}（不用它会怎样）、\textbf{影响方向}与\textbf{验证方式}，
凡无法给出验证方式的假设一律删除。

\begin{table}[htbp]
  \centering
  \caption{模型假设及其必要性、影响方向与验证方式}
  \label{tab:assumptions}
  \footnotesize
  \setlength{\tabcolsep}{4pt}
  \begin{tabularx}{\linewidth}{C{0.7cm} L{2.9cm} L{3.4cm} L{2.9cm} L{1.7cm} Y}
    \toprule
    编号 & 假设 & 必要性（不用会怎样） & 现实/题目依据 & 影响方向 & 验证或灵敏度方式 \\
    \midrule
    A1 & 不允许向外部电网售电
       & 目标函数将多出一项售电收益，盈余有出路，弃光恒为 0
       & 题面未给出上网电价；就近消纳"优先自用"导向
       & 使费用\textbf{偏高}，结果保守
       & 引入上网电价 \(\theta\) 扫描，观察最优策略与费用变化 \\
    \addlinespace
    A2 & 充电效率与放电效率各取 0.9（往返 0.81）
       & 效率决定套利是否有利可图，是收益的核心参数
       & 题面给出"充放电效率为 90\%"，\textbf{未区分充、放方向}
       & 效率降低 \(\Rightarrow\) 套利收益降低
       & 0.85、0.90、0.95 三档 + "单侧 90\%"与"往返 81\%"双口径对照 \\
    \addlinespace
    A3 & 同一时段不同时充放电
       & 若允许同时充放，须引入 0--1 变量，模型由 LP 变为 MILP，规模与求解时间显著上升
       & 效率 \(<1\) 时同时充放严格劣于只做净动作（可证，见命题~\ref{prop:no-simul}）
       & 不影响最优目标值
       & 数值断言 \(\min(c_t,d_t)=0\) 全部成立（审计 A5） \\
    \addlinespace
    A4 & 忽略自放电与循环衰减
       & 否则目标函数须加入与循环次数相关的衰减成本项
       & 题面未给出衰减参数
       & 高估长期收益
       & 加入循环成本 \(c_{cyc}\) 做灵敏度，给出结论反转点 \\
    \addlinespace
    A5 & 储电量跨日连续递推（不逐日清零）
       & 若每日清零，则与"2025-01-01 \(0{:}00\) 为 6000~kWh"的初值条件矛盾，且丢弃跨日套利空间
       & 附录 1 给出的初值 + 储能运行常识
       & 影响全窗口费用
       & 与"逐日独立（每日 SOC 清零边界）"变体对照（消融层级 B2） \\
    \addlinespace
    A6 & 结果窗口末端施加终端条件 \(S_T\ge \underline S\)（并给出日周期变体）
       & 否则模型会在窗口末日"掏空"储能以降低当期费用，形成短视解
       & 储能运营常识；问题 1 题面直接要求 \(S_0=S_{144}\)
       & 影响末期费用
       & 终端处理四组对照：自由、\(S_T\ge S_0\)、日周期、水价值终端（§9.4 层级 B5） \\
    \addlinespace
    A7 & 调整决策"后覆盖前"，调整费用的计算基准始终为 \(0{:}00\) 的计划购电量
       & 若基准随上一次调整漂移，调整费与违约费口径不可复现
       & 题面"计划购电量高于调整购电量的部分"表述
       & 影响调整相关费用
       & 两套口径（相对 \(0{:}00\) 计划、相对上次调整）对照并各自命名 \\
    \addlinespace
    A8 & 2025 年 1 月作为储电量递推的预热期（warm-up）
       & 结果窗口自 2025-02-01 起，但初值自 2025-01-01 起算，若不作预热声明则初期储电量无来源
       & 附件 2 覆盖 2025-01-01 至 12-31，而结果文件要求 2025-02-01 起
       & 影响初期储电量
       & 论文中明示 1 月不计入结果；对 1 月处理方式做对照 \\
    \bottomrule
  \end{tabularx}
\end{table}

\begin{guideblock}
两条删除规则：① "建议保留一位小数""数据已清洗"属数据处理，不是模型假设；
② 未被任何公式使用的假设属\textbf{装饰性假设}，删掉反而更专业。
表格的"验证方式"列\textbf{必须}填出可执行的档位（如"0.85/0.90/0.95 三档"），
填不出验证方式的假设，直接删除。
\end{guideblock}

\subsection{信息结构声明}

\begin{guideblock}
这是本题\textbf{最大的争议点}，必须在假设节显式声明，并说明"计划依据"与"结算依据"
是两套数据。三种解读的后果差异极大，选错会让整篇论文失去意义。
\end{guideblock}

问题 2--4 中"每天 \(0{:}00\) 制定计划"这一条件，存在三种可能的解读：

\begin{table}[htbp]
  \centering
  \caption{信息结构的三种解读及其后果}
  \label{tab:infostruct}
  \footnotesize
  \begin{tabularx}{\linewidth}{C{1.0cm} L{4.6cm} L{4.6cm} Y}
    \toprule
    代号 & 解读 & 后果 & 本文的使用方式 \\
    \midrule
    R1 & 完美信息：\(0{:}00\) 即知当日负载与光伏真值
       & 紧急购电恒为 0，问题 2 退化为问题 1 的重复，题面表 3 失去意义
       & \textbf{不可作为主方案}；用作"预测误差成本"的下界对照 \\
    \addlinespace
    R2 & 典型日预报：以附件 1 的日均形状作为预报
       & 预报被设定得很差，紧急购电量最多，结果偏悲观
       & 仅作悲观对照；附件 1 是实际值年均而非预报，若用作 R2 须额外论证并做泄漏检验，且不得作为主方案 \\
    \addlinespace
    R3 & 日前预报：问题 2/3 使用附件 3 的 \(0{:}00\) 预报（问题 2 无预报数据时按题面采用当日负载与光伏的实际可得信息结构另行声明）
       & 紧急购电量中等，最符合工程实际与题面"制定计划"的语义
       & \textbf{推荐主方案}；\note{如最终选用 R2 或自定义口径，务必在此处改写并说明理由} \\
    \bottomrule
  \end{tabularx}
\end{table}

本文采用\textbf{信息结构 R3}：计划依据为预报值（问题 3 中为附件 3 的 \(0{:}00\) 预报），
结算依据为附件 2 的当日实际值，二者是\textbf{两套数据}。
无论采用哪种解读，本文都做到三件事：① 在本节显式声明；
② 说明计划依据与结算依据的分离方式；③ 用灵敏度检验结论对预测误差水平的稳健性
（预报误差 \(\times 0.5/\times 1/\times 2\) 三档）。

\subsection{符号说明}

\begin{guideblock}
公式中的每一个符号都必须在本表出现，避免"天降符号"。
\textbf{对偶变量也是模型的组成部分，必须登记}，否则第五章的阈值律无法被读懂。
单位要写清楚，特别注意 kW（功率）与 kWh（电量）不得混用。
\end{guideblock}

\begin{table}[htbp]
  \centering
  \caption{符号说明（集合、参数、决策变量与对偶变量）}
  \label{tab:notation}
  \footnotesize
  \begin{tabularx}{\linewidth}{C{1.2cm} C{2.6cm} L{5.2cm} C{1.6cm} Y}
    \toprule
    类别 & 符号 & 含义 & 单位 & 来源/取值 \\
    \midrule
    集合 & \(t\in\mathcal T\) & 时段（10~min 粒度） & --- & \(|\mathcal T|=144\) \\
    集合 & \(d\in\mathcal D\) & 日期 & --- & 334 天（2025-02-01 至 12-31） \\
    集合 & \(m\in\mathcal M\) & 预报发布/决策时刻 & --- & \(\{0,6,12,18\}\) 时 \\
    集合 & \(\omega\in\Omega\) & 场景 & --- & \(|\Omega|=N\)，由 §5.2 确定 \\
    \addlinespace
    参数 & \(L_{d,t}\) & 小区负载（功率） & kW & 附件 2 \\
    参数 & \(PV_{d,t}\) & 光伏实际功率 & kW & 附件 2 \\
    参数 & \(\widehat{PV}^{(m)}_{d,h}\) & 光伏预报功率 & kW & 附件 3 \\
    参数 & \(p_{d,t}\) & 电价 & 元/kWh & 附件 1（固定价）/ 附件 4（波动价） \\
    参数 & \(\Delta t\) & 时间步长 & h & \(1/6\) \\
    参数 & \(\eta_c,\eta_d\) & 充电效率、放电效率 & --- & \(0.9/0.9\)（假设 A2） \\
    参数 & \(\overline S,\underline S\) & 储电量上下限 & kWh & \(10800/1200\) \\
    参数 & \(\overline S^{\max}\) & 储能设备最大容量 & kWh & 12000 \\
    参数 & \(\overline P\) & 最大充放电功率 & kW & 5000 \\
    参数 & \(S^{init}\) & 初始储电量 & kWh & 6000 @ 2025-01-01 \(0{:}00\) \\
    参数 & \(\kappa\) & 紧急购电电价倍数 & --- & 5 \\
    参数 & \(\kappa^{+},\kappa^{-}\) & 调整超额 / 违约电价倍数 & --- & \(1.5/0.5\) \\
    参数 & \(\beta\) & CVaR 置信水平 & --- & 0.95（预先锁定） \\
    \addlinespace
    决策 & \(E^{buy}_{d,t}\) & 计划购电量 & kWh & --- \\
    决策 & \(E^{adj}_{d,t}\) & 调整购电量 & kWh & --- \\
    决策 & \(g_{d,t}\) & 紧急购电量 & kWh & --- \\
    决策 & \(c_t,d_t\) & 充电量、放电量 & kWh & \(\le \overline P\Delta t\) \\
    决策 & \(s_t\) & 弃光电量 & kWh & \(\ge 0\) \\
    决策 & \(S_t\) & 时段末储电量 & kWh & \(\in[\underline S,\overline S]\) \\
    \addlinespace
    \textbf{对偶} & \(\mu_t\) & 储电的边际价值（水价值） & 元/kWh & LP 对偶 \\
    \textbf{对偶} & \(\lambda^{\pm}_t\) & 储电量上下界的影子价格 & 元/kWh & LP 对偶 \\
    \textbf{对偶} & \(z^{c}_t,z^{d}_t\) & 充/放互补松弛变量 & 元/kWh & \(z^c_t=p_t-\eta_c\mu_{t}\)；\(z^d_t=-p_t+\mu_{t}/\eta_d\) \\
    \bottomrule
  \end{tabularx}
\end{table}

\begin{guideblock}
符号表自检：① 每个在正文公式里出现过的符号是否都在表中？
② 对偶变量 \(\mu_t\)、\(\lambda^{\pm}_t\)、\(z^{c}_t\)、\(z^{d}_t\) 是否已登记？（第五章要用）
③ 全文统一称"微网"，不混用"微电网"；统一使用题面名词
"储能设备""购电量、计划购电量、调整购电量、紧急购电量""外网"；
④ 首次出现的外文术语用"中文（英文）"形式。
\end{guideblock}
